% ====================================================================

% ====================================================================
\documentclass[article,paper=a4paper]{jlreq}

% jpreportパッケージの読み込み
\usepackage{jpreport}

% ページレイアウトの設定
\usepackage{geometry}
\usepackage{enumitem}
\geometry{
    top=25mm,
    bottom=25mm,
    left=25mm,
    right=25mm
}

% ハイパーリンクの設定
\usepackage[unicode]{hyperref}
\usepackage{bookmark}
\usepackage[nameinlink]{cleveref}

\hypersetup{
    colorlinks=true,
    linkcolor=black,
    urlcolor=blue,
    citecolor=blue,
    pdftitle=,
    pdfauthor=
}

% 数学演算子の定義


% cleveref用の日本語設定
\crefname{thm}{定理}{定理}
\crefname{lem}{補題}{補題}
\crefname{cor}{系}{系}
\crefname{ex}{例}{例}
\crefname{dfn}{定義}{定義}
\crefname{prop}{命題}{命題}
\crefname{prob}{問題}{問題}

% 色彩の定義
\definecolor{yukari}{RGB}{196, 163, 191}
% 文書情報
\title{解析学1}
\author{akari0koutya}
\date{\today}

\begin{document}

% --- タイトル ------------------------------------------------------
\maketitle
\tableofcontents
%% 本文の内容をここに記述
% ====================================================================

\section{ガイダンス (2026/04/14)}
\subsection{注意点}
\begin{itemize}
    \item 5回以内の欠席であれば,期末試験を受けれる。期末試験は出席回数にカウントしない。
    \item 2回の遅刻で1回の欠席とみなす。
    \item 紙面での欠席確認を行うことがある。
    \item レポートの提出場所は数理のレポートボックス。
    \item 未完成のレポートは未提出扱い。
    \item 成績評価は $\text{期末試験}:\text{レポート} = 3:1$
\end{itemize}

\subsection{アブストラクト}

\begin{tcolorbox}[colback=white,colframe=yukari,title=公理的確率の存在]
$\Omega : \text{non empty set}$ \par
集合関数 $P$ が次を満たすとき、 $P$ を $\Omega$ 上の（\textbf{公理的}）\textbf{確率}と呼ぶ :
\begin{itemize}
    \item $P(\Omega) = 1$
    \item 任意の $A \subseteq \Omega$ に対して、$P(A) \geq 0$ 
    \item 非交な集合列 $A_1, A_2, \ldots$ に対して、 $P\left( \bigcup_{i=1}^{\infty} A_i \right) = \sum_{i=1}^{\infty} P(A_i)$
\end{itemize}
\end{tcolorbox}

\begin{tcolorbox}[colback=white,colframe=yukari,title=リーマン積分]
$f \colon [a,b] \to \mathbb{R}$ を区間 $[a,b]$ 上の関数とする。 $f$ が \textbf{リーマン積分可能} であるとは、次の条件を満たすことをいう :
\begin{itemize}
    \item 任意の $\epsilon > 0$ に対して、区間 $[a,b]$ を分割する分割 $P$ が存在し、任意の $P$ に対して、上和 $U(f,P)$ と下和 $L(f,P)$ の差が $\epsilon$ より小さいこと。
\end{itemize}
\end{tcolorbox}

\begin{tcolorbox}[colback=white,colframe=yukari,title=リーマン積分の欠点1]
    リーマン積分可能な関数を特徴づけることができない。
    \begin{prop}
        $f$ が区間 $[a,b]$ でリーマン積分可能ならば、 $f$ は区間 $[a,b]$ 上で有界。
    \end{prop}
    これは逆が成り立たない。例えば、 $f(x) = \begin{cases} 1 & x \in \mathbb{Q} \\ 0 & x \in \mathbb{R} \setminus \mathbb{Q} \end{cases}$ は区間 $[0,1]$ 上で有界だが、リーマン積分可能ではない。
\end{tcolorbox}

\begin{tcolorbox}[colback=white,colframe=yukari,title=リーマン積分の欠点2]
    リーマン積分は、項別積分に関して問題がある。\par
    例えば、 $f_n(x) = \begin{cases} 1 & x \in [0,1/n] \\ 0 & x \in (1/n,1] \end{cases}$ とすると、各 $f_n$ は区間 $[0,1]$ 上でリーマン積分可能であり、 $\int_0^1 f_n(x) dx = 1/n$ となる。しかし、 $\lim_{n \to \infty} f_n(x) = 0$ であるため、 $\int_0^1 \lim_{n \to \infty} f_n(x) dx = 0$ となる。つまり、 $\lim_{n \to \infty} \int_0^1 f_n(x) dx \neq \int_0^1 \lim_{n \to \infty} f_n(x) dx$。
\end{tcolorbox}

これらの問題を解決するために、ルベーグ積分が導入される。

\subsection{教科書}
\begin{itemize}
    \item 吉田洋一, ルベグ積分入門, ちくま学芸文庫.
\end{itemize}
上のアブストラクトは、「ルベグ積分入門」の1章の内容を簡単にまとめたものである。
\end{document}